The MCU has several \peripheral{GPIOx} peripherals, each of which controls a set of general purpose input/output (GPIO) pins. The VestaRV implementation supports GPIO ports with 8, 16, or 32 pins. Each GPIO peripheral (typically called a port) constitutes a parallel port, allowing all pins in the port to be updated simultaneously. Each GPIO pin may also be configured and updated individually, allowing for design flexibility. Each GPIO pin is referred to as Px.y, where x refers to the port number and y refers to the corresponding bit in the GPIO registers that controls the pin. For example, P2.5 refers to port 2, bit 5.

Every GPIO pin on the MCU can be in one of two modes: primary/GPIO mode, or secondary/peripheral mode. Section \ref{s:pinsConfig} details the secondary function of each GPIO pin.



\subsection{Primary/GPIO Mode Functionality}

In primary/GPIO mode, the GPIO pin state is controlled by the GPIO registers \register{PxDIR}, \register{PxOUT}, and \register{PxREN}. The \register{PxDIR} register determines if the pin is in input or output mode, the \register{PxOUT} register determines if the pin outputs a logic high or logic low level when the pin is in output mode, and the \register{PxREN} register enables or disables the internal pullup/pulldown resistor. Table \ref{t:gpio-config} shows the various states available to every GPIO pin.

\begin{table}[htb]
	\centering
	\begin{tabular}{ c c c c }
	\caption{GPIO Configurations} \label{t:gpio-config}
	\textbf{\register{PxDIR[y]}} & \textbf{\register{PxOUT[y]}} & \textbf{\register{PxREN[y]}} & \textbf{Px.y State} \\ \hline
	0 & - & 0 & High impedance input, resistor disabled \\
	\rowcolor{tablehighlightcolor}
	0 & - & 1 & Pullup/pulldown resistor enabled \\
	1 & 0 & - & Output, strongly driven low \\
	\rowcolor{tablehighlightcolor}
	1 & 1 & - & Output, strongly driven high \\
	\hline
	\end{tabular}
\end{table}

The \register{PxIN} register always reads the digital state of the pin, regardless of whether it is in input or output mode, and regardless of whether it is in primary/GPIO mode or secondary/peripheral mode. The pin Px.y is at logic low level if \register{PxIN[y]} is 0, and is at logic high level if it is 1. Note that \register{PxIN} is latched on the falling edge of the memory enable signal to ensure stable readings during memory access.

\subsection{Register Access Convenience Features}

The GPIO peripheral provides convenient set, clear, and toggle registers for efficient bit manipulation without read-modify-write operations:

\begin{itemize}
\item \register{PxOUTS}: Writing 1 to a bit sets the corresponding output pin high. Writing 0 has no effect. Reading returns the current \register{PxOUT} value.
\item \register{PxOUTC}: Writing 1 to a bit clears the corresponding output pin low. Writing 0 has no effect. Reading returns the inverted \register{PxOUT} value.
\item \register{PxOUTT}: Writing 1 to a bit toggles the corresponding output pin. Writing 0 has no effect. Reading returns the current \register{PxOUT} value.
\end{itemize}

These registers eliminate race conditions in interrupt-driven or multi-threaded code where multiple code paths may modify GPIO outputs.


\subsection{Secondary/Peripheral Mode Functionality}

Most GPIO pins have a secondary/peripheral function which is controlled by one of the peripherals in the MCU. The mode of the GPIO is determined by the \register{PxSEL} register. If \register{PxSEL[y]} is 0, Px.y is in primary/GPIO mode, and if it is 1, Px.y is in secondary/peripheral mode.

When in peripheral mode, the peripheral seizes control over the pin's \register{DIR}, \register{OUT}, and \register{REN} controls, overriding \register{PxDIR[y]}, \register{PxOUT[y]}, and \register{PxREN[y]}. For example, when the \pin{SCK0} pin is in peripheral mode, the \peripheral{SPI0} peripheral controls the state of the pin, and the corresponding \register{PxDIR[y]}, \register{PxOUT[y]}, and \register{PxREN[y]} registers have no effect until \register{PxSEL[y]} returns to 0.

Some pins with secondary/peripheral functions may only be used as inputs by the peripheral. For example, timer capture pins like \pin{T0CAP0} must be manually configured as inputs in \register{PxDIR} before use, even when \register{PxSEL[y]} is set to enable the secondary function.



\subsection{Pin Interrupts}

The VestaRV GPIO implementation provides individual interrupt outputs for each GPIO pin. Each pin can generate an interrupt independently, allowing for flexible interrupt handling through the system's interrupt vector table.

To enable interrupts on a pin:
\begin{enumerate}
\item Unmask the corresponding GPIO pin interrupt in the system interrupt controller
\item Set \register{PxIE[y]} to 1 for the desired pin Px.y
\item Configure the interrupt edge using \register{PxIES[y]}: 0 for low-to-high (rising edge), 1 for high-to-low (falling edge)
\end{enumerate}

When the pin Px.y experiences the edge selected by \register{PxIES[y]} while \register{PxIE[y]} is enabled, the interrupt flag \register{PxIF[y]} is immediately set to 1 and the corresponding pin-specific ISR is executed. The interrupt flag must be cleared by writing 1 to \register{PxIF[y]} before the ISR returns, otherwise the interrupt will re-trigger.

\textbf{Important:} Modifying \register{PxIES[y]} while \register{PxIE[y]} is enabled may immediately generate a spurious interrupt. Always disable the pin interrupt (\register{PxIE[y]} = 0) before changing \register{PxIES[y]}.

The \register{PxIF} register is also latched on the falling edge of the memory enable signal to ensure stable readings during interrupt service routines.
